\section{Introduction}
\label{sec:introduction}

Image denoising is a fundamental problem in low-level computer vision, aiming to recover a clean image from its noisy observation. With the advent of deep learning, CNN-based approaches have largely superseded traditional methods such as BM3D and NLM, achieving state-of-the-art results across a variety of noise models and intensity levels~\cite{dncnn}.

This work presents a systematic comparison of multiple deep learning architectures for image restoration. Rather than proposing a single novel method, we focus on implementing a \emph{unified experimental framework} that enables fair, reproducible comparisons across fundamentally different architectural paradigms: encoder-decoder networks with skip connections (UNet family), purely sequential residual networks (DnCNN), attention-augmented architectures (RIDNet, Attention UNet), activation-free networks (NAFNet), and generative adversarial networks (Pix2Pix).

All models are evaluated under identical conditions: the same dataset (DIV2K~\cite{div2k}), the same degradation model (additive Gaussian noise with $\sigma = 100$, a deliberately severe setting), and the same evaluation protocol (sliding-window inference with Hann blending on full-resolution images).

\subsection{Contributions}

The main contributions of this work are:
\begin{itemize}
    \item A modular, extensible framework covering the complete pipeline from degradation generation through training to evaluation;
    \item A fair comparison of eight architectures under controlled conditions, with extensive ablation studies;
    \item Analysis of critical design choices: residual learning vs.\ direct prediction, bilinear vs.\ transposed convolution upsampling, batch normalization vs.\ group normalization, and multi-component loss functions.
\end{itemize}

\subsection{Work Division}

The project was developed collaboratively:
\begin{itemize}
    \item \textbf{Giuseppe Bellamacina} --- UNet, UNet Residual, Attention UNet, and the shared infrastructure (degradations, training pipeline, evaluation, utilities);
    \item \textbf{Salvo Iurato} --- DnCNN and RIDNet;
    \item \textbf{Mattia Campanella} --- Denoising Autoencoder and NAFNet scaling study ($\sigma\!=\!50$);
    \item \textbf{Daniele Barbagallo} --- Pix2Pix (conditional GAN) and NAFNet-V2 ($\sigma\!=\!100$).
\end{itemize}

\subsection{Paper Organization}

Section~\ref{sec:degradations} describes the degradation models. Section~\ref{sec:infrastructure} details the shared training infrastructure. Section~\ref{sec:models} presents the model architectures. Section~\ref{sec:training} discusses training procedures and hyperparameters. Section~\ref{sec:evaluation} defines the evaluation methodology. Section~\ref{sec:results} reports experimental results with ablation studies and visual comparisons. Section~\ref{sec:conclusion} concludes the paper.
